% Motor: XeLaTeX. Sesión 28 - material docente UTP.
% Estructura docente de la Sesión 27.
\documentclass[aspectratio=169,10pt,c]{beamer}
% Reutiliza sin cambios la identidad visual UTP validada en la Sesion 08.
% Reutiliza sin cambios la identidad visual UTP validada en la Sesion 05.
% La segunda ruta permite compartir los assets comunes sin duplicar binarios.
% !TEX root = slides.tex
% Preambulo comun: estilo UTP, colores, fuentes, plantillas de beamer y
% comandos de apoyo didactico. Se carga desde slides.tex tras \documentclass.

% La opción "c" de beamer reparte el espacio libre 1:1.5 (arriba:abajo), por lo
% que el contenido queda por encima del centro. Igualamos ambos rellenos.
\makeatletter
\define@key{beamerframe}{c}[true]{%
  \beamer@frametopskip=0pt plus 1fill\relax%
  \beamer@framebottomskip=0pt plus 1fill\relax%
  \beamer@frametopskipautobreak=0pt plus .5\paperheight\relax%
  \beamer@framebottomskipautobreak=0pt plus .5\paperheight\relax%
}
\makeatother

\usepackage{graphicx}
\usepackage{tikz}
\usepackage{xcolor}
\usepackage{iftex}

\ifPDFTeX
  \usepackage[utf8]{inputenc}
  \usepackage[T1]{fontenc}
  \usepackage{lmodern}          % tipografia escalable para el respaldo con pdflatex
  \renewcommand{\familydefault}{\sfdefault}
\else
  \usepackage{fontspec}
  \IfFontExistsTF{Arial}
    {\setsansfont{Arial}[BoldFont={Arial Bold}]}
    {\setsansfont{Helvetica}[BoldFont={Helvetica Bold}]}
\fi

\usepackage[spanish,es-nodecimaldot]{babel}
\usepackage{microtype}
\usepackage{booktabs}
\usepackage{array}
\usetikzlibrary{arrows.meta,positioning,fit,calc,shapes.geometric}

\graphicspath{{assets/}}

% ---------- Identidad visual UTP ----------
\definecolor{utpRojo}{HTML}{C42B35}
\definecolor{utpNegro}{HTML}{111111}
\definecolor{utpGris}{HTML}{5F5F5F}
\definecolor{utpGrisMedio}{HTML}{9A9A9A}
\definecolor{utpGrisClaro}{HTML}{F1F1F1}

\usetheme{default}
\setbeamercolor{normal text}{fg=utpNegro,bg=white}
\setbeamercolor{structure}{fg=utpRojo}
\setbeamercolor{frametitle}{fg=utpNegro,bg=white}
\setbeamercolor{title}{fg=utpNegro}
\setbeamercolor{subtitle}{fg=utpNegro}
\setbeamercolor{author}{fg=utpNegro}
\setbeamercolor{date}{fg=utpNegro}
\setbeamercolor{block title}{fg=white,bg=utpNegro}
\setbeamercolor{block body}{fg=utpNegro,bg=utpGrisClaro}
\setbeamercolor{alertblock title}{fg=white,bg=utpRojo}
\setbeamercolor{alertblock body}{fg=utpNegro,bg=utpRojo!7}
\setbeamercolor{exampleblock title}{fg=white,bg=utpGris}
\setbeamercolor{exampleblock body}{fg=utpNegro,bg=utpGris!8}
\setbeamercolor{terminal}{fg=white,bg=utpNegro}

\setbeamerfont{title}{series=\bfseries,size=\fontsize{19.5}{22.5}\selectfont}
\setbeamerfont{subtitle}{size=\fontsize{11.5}{13.5}\selectfont}
\setbeamerfont{date}{series=\bfseries,size=\fontsize{11.5}{13.5}\selectfont}
\setbeamerfont{author}{size=\fontsize{10.5}{12.5}\selectfont}
\setbeamerfont{frametitle}{series=\bfseries,size=\large}
\setbeamerfont{block title}{series=\bfseries,size=\small}
\setbeamerfont{block body}{size=\normalsize}

\setbeamertemplate{navigation symbols}{}
\setbeamertemplate{headline}{}
\setbeamertemplate{blocks}[default]
\setbeamertemplate{itemize item}{\color{utpRojo}\small\raise.15ex\hbox{$\blacktriangleright$}}
\setbeamertemplate{itemize subitem}{\color{utpGris}\scriptsize$\bullet$}
\setbeamertemplate{enumerate item}{\color{utpRojo}\insertenumlabel.}
\setbeamertemplate{background canvas}{\color{white}\rule{\paperwidth}{\paperheight}}

% Título de diapositiva: negrita negra a la izquierda, logo UTP a la derecha
\setbeamertemplate{frametitle}{%
  \begin{tikzpicture}[remember picture,overlay]
    \node[anchor=north east,inner sep=0pt]
      at ([xshift=-7mm,yshift=-5.5mm]current page.north east)
      {\includegraphics[width=17mm]{utp-logo.png}};
  \end{tikzpicture}%
  \vspace{1.2mm}%
  \begin{minipage}[t]{.80\textwidth}\raggedright\insertframetitle\end{minipage}\par
  \vspace{1.1mm}{\color{utpRojo}\rule{\textwidth}{1.1pt}}
}

% Solo numeración discreta; la plantilla UTP no lleva pie de página
\setbeamertemplate{footline}{%
  \hfill{\color{utpGrisMedio}\scriptsize\insertframenumber}\hspace{6mm}\vspace{2.5mm}%
}

% Portada UTP
\setbeamertemplate{title page}{%
  \begin{tikzpicture}[remember picture,overlay]
    \node[anchor=center,inner sep=0pt] at (current page.center) {%
      \includegraphics[width=\paperwidth,height=\paperheight]{utp-cover-bg.jpeg}%
    };
    \node[anchor=north,align=center,text=utpNegro,inner sep=0pt,
      text width=145mm,font=\usebeamerfont{title}]
      at ([yshift=-25.2mm]current page.north) {\inserttitle};
    \node[anchor=north,align=center,text=utpNegro,inner sep=0pt,
      text width=142mm,font=\usebeamerfont{subtitle}]
      at ([yshift=-38.2mm]current page.north) {\insertsubtitle};
    \node[anchor=north,align=center,text=utpNegro,inner sep=0pt,
      font=\usebeamerfont{date}]
      at ([yshift=-50.4mm]current page.north) {\insertdate};
    \node[anchor=north,align=center,text=utpNegro,inner sep=0pt,
      font=\usebeamerfont{author}]
      at ([yshift=-58.4mm]current page.north) {\insertauthor};
  \end{tikzpicture}%
}

% ---------- Comandos de apoyo didáctico ----------
\newcommand{\concepto}[2]{%
  \begin{block}{#1}\large #2\end{block}}
\newcommand{\pregunta}[1]{%
  \begin{alertblock}{Decide y explica}#1\end{alertblock}}
\newcommand{\analogia}[1]{%
  \begin{exampleblock}{En palabras simples}\small #1\end{exampleblock}}
\newcommand{\enlarealidad}[1]{%
  \begin{block}{Cómo se ve en la realidad}\small #1\end{block}}
\newcommand{\ojo}[1]{%
  \begin{alertblock}{Ojo con esto}\small #1\end{alertblock}}
\newcommand{\apoyo}{{\normalfont\footnotesize\color{utpGris}~· material de apoyo}}
\newcommand{\marca}[1]{\tikz[baseline=(m.base)]{%
  \node[rounded corners=1.5pt,fill=utpGrisClaro,draw=utpGrisMedio,inner xsep=3pt,inner ysep=1.5pt,
        font=\footnotesize] (m) {#1};}}
\newenvironment{terminal}{%
  \begin{beamercolorbox}[wd=\linewidth,sep=5pt,rounded=true]{terminal}%
  \ttfamily\small\ignorespaces}{%
  \end{beamercolorbox}}

\tikzset{
  caja/.style={draw=utpNegro, very thick, rounded corners=1.5pt, fill=white,
    align=center, minimum height=0.72cm, inner sep=5pt, font=\small},
  cajaazul/.style={caja, draw=utpGris, fill=white},
  cajaverde/.style={caja, draw=utpNegro, fill=utpGrisClaro},
  cajanaranja/.style={caja, draw=utpRojo!75, fill=utpRojo!7},
  flecha/.style={-{Stealth[length=2.3mm]}, very thick, draw=utpNegro},
  flechavuelta/.style={{Stealth[length=2.3mm]}-{Stealth[length=2.3mm]}, very thick, draw=utpNegro}
}

\graphicspath{{../../Sesion05/slides/assets/}}



\graphicspath{{assets/}{../../Sesion05/slides/assets/}}
\usepackage{colortbl}
\usepackage{listings}
\lstset{basicstyle=\ttfamily\footnotesize,backgroundcolor=\color{utpGrisClaro},
  frame=none,aboveskip=4pt,belowskip=4pt,showstringspaces=false,
  columns=fullflexible,keepspaces=true,breaklines=true,
  keywordstyle=\color{utpRojo}\bfseries,
  morekeywords={FROM,WORKDIR,COPY,RUN,EXPOSE,CMD},
  commentstyle=\color{utpGris}\itshape,morecomment=[l]{\#},
  xleftmargin=5pt,xrightmargin=5pt,framesep=5pt}
% El código de Python: solo comentarios y cadenas con color.
\lstdefinestyle{python}{basicstyle=\ttfamily\scriptsize,keywords={},
  stringstyle=\color{utpRojo},morestring=[b]",commentstyle=\color{utpGris}\itshape}
% Un compose.yaml: las claves en rojo.
\lstdefinestyle{yaml}{basicstyle=\ttfamily\scriptsize,
  keywords={services,api,web,build,image,ports},keywordstyle=\color{utpRojo}\bfseries}
% Una terminal: fondo oscuro, sin colorear palabras.
\lstdefinestyle{terminal}{basicstyle=\ttfamily\scriptsize\color{white},backgroundcolor=\color{utpNegro},
  keywords={},morecomment=[l]{\#},commentstyle=\color{utpGrisMedio},escapeinside={(*}{*)},upquote=true}
\usepackage{upquote}
\usepackage{pifont}
\definecolor{okVerde}{HTML}{2E7D32}
\newcommand{\bien}{\textcolor{okVerde}{\ding{51}}}
\newcommand{\mal}{\textcolor{utpRojo}{\ding{55}}}
% Un número en un círculo, para señalar líneas de una terminal.
\newcommand{\nro}[1]{\tikz[baseline=(m.base)]\node[circle,fill=utpRojo,text=white,inner sep=1pt,font=\tiny\bfseries\sffamily](m){#1};}

\title{Integración de Servicios Cloud GCP}
\subtitle{Docker: tu propia imagen, puertos, variables y Compose}
\date{Semana 14 - Sesión 28}
\author{Rel Guzman}

\newcommand{\separador}[2]{%
\begin{frame}[plain]
\vfill
\begin{center}
{\fontsize{29}{34}\selectfont\bfseries #1}\par
\vspace{2mm}
{\small #2}\par
\vspace{4mm}
{\color{utpRojo}\rule{8.5cm}{1.2pt}}
\end{center}
\vfill
\end{frame}}

% Íconos simples para los diagramas.
\tikzset{
 caja/.style={draw=utpRojo,thick,rounded corners=4pt,align=center},
 vm/.style={draw=utpRojo,thick,rounded corners=3pt,fill=white,minimum width=1.6cm,minimum height=0.75cm,align=center,font=\scriptsize},
 zona/.style={draw=utpGrisMedio,dashed,rounded corners=4pt},
 flecha/.style={-latex,thick},
 capa/.style={draw=utpRojo,thick,rounded corners=2pt,minimum width=5.6cm,minimum height=0.48cm,font=\scriptsize\ttfamily,anchor=west},
}

\begin{document}

{\setbeamertemplate{footline}{}\maketitle}

% 2
\begin{frame}{Revisión: lo que vimos en la Sesión 27}
\begin{columns}[T,onlytextwidth]
\column{.46\textwidth}
\begin{block}{Desplegar y dependencias}
\small Desplegar es dejar una app funcionando en un servidor. Una dependencia es lo que la app necesita y no es su código: Python, Flask.
\end{block}
\begin{block}{Imagen y contenedor}
\small La imagen es la app con sus dependencias, lista para copiar. Un contenedor es una copia de la imagen en marcha.
\end{block}
\begin{block}{Ciclo de vida}
\small \texttt{docker run} lo crea y lo arranca, \texttt{docker stop} lo detiene y \texttt{docker rm} lo borra.
\end{block}
\column{.50\textwidth}
\begin{alertblock}{Responde}
\small
1. Diego copió \texttt{api.py} a su laptop y salió \texttt{ModuleNotFoundError}. ¿Qué le faltaba?\par\medskip
2. Lucía le pasa a Diego la \textbf{imagen} de la API. ¿Qué necesita Diego para correrla?\par\medskip
3. Detuviste \texttt{web} con \texttt{docker stop}. ¿Con qué comando lo ves todavía? ¿Cómo lo borras?
\end{alertblock}
\end{columns}
\end{frame}

% 3
\begin{frame}{Revisión: la pregunta de hoy}
\begin{center}
\begin{tikzpicture}[font=\small]
  \node[caja,fill=utpGrisClaro,text width=2.5cm,minimum height=1.2cm] (cod) at (0,0) {\textbf{La API}\\\scriptsize\texttt{api.py} y\\\texttt{requirements.txt}};
  \node[caja,fill=utpRojo!30,text width=2.0cm,minimum height=1.2cm] (img) at (3.9,0) {\textbf{Imagen}\\\footnotesize\texttt{cafe-api}};
  \node[caja,draw=okVerde,fill=okVerde!12,text width=2.2cm,minimum height=1.2cm] (k) at (7.8,0) {\textbf{Contenedor}\\\footnotesize la API en marcha};
  \node[caja,dashed,text width=2.3cm,minimum height=1.2cm] (nav) at (11.45,0) {\textbf{Navegador}\\\footnotesize\texttt{localhost:8080}};
  \draw[flecha,dashed] (cod) -- node[above,font=\bfseries]{?} (img);
  \draw[flecha] (img) -- node[above,font=\scriptsize\ttfamily]{docker run} (k);
  \draw[flecha,dashed] (nav) -- node[above,font=\bfseries]{?} (k);
\end{tikzpicture}
\end{center}
\vspace{3mm}
\begin{columns}[T,onlytextwidth]
\column{.32\textwidth}
\begin{block}{1. Dockerfile}
\small ¿Cómo se arma una imagen propia con la API de Café Nube?
\end{block}
\column{.32\textwidth}
\begin{block}{2. Puertos y variables}
\small ¿Cómo le llega un pedido al contenedor? ¿Cómo se configura sin cambiar la imagen?
\end{block}
\column{.32\textwidth}
\begin{block}{3. Compose}
\small Si la app tiene varias piezas, ¿cómo se levantan todas juntas?
\end{block}
\end{columns}
\end{frame}

% 4
\begin{frame}[c]
\begin{block}{Logro de la sesión}
Al finalizar, podrás explicar con tus palabras qué hace cada línea del \textbf{Dockerfile} de la API de Café Nube, cómo se \textbf{publica un puerto} y se pasa una \textbf{variable de entorno} a un contenedor, por qué importa el orden de las \textbf{capas}, y cómo \textbf{Docker Compose} levanta y baja una app de varias piezas con un comando.
\end{block}
\end{frame}

% 5
\begin{frame}{Saberes previos}
\begin{itemize}\setlength\itemsep{14pt}
  \item En la Sesión 27, ¿qué tuviste que instalar en tu computadora antes de encender la API de Café Nube?
  \item Cuando encendiste la API por segunda vez en el 8080, ¿qué pasó? ¿Cómo lo resolviste sin cambiar el código?
  \item ¿Qué hace \texttt{docker run hello-world} si la imagen no está en tu computadora?
\end{itemize}
\end{frame}

% 6
\begin{frame}[c]{Utilidad de esta sesión}
\begin{center}
\begin{minipage}{.84\textwidth}
{\large Hasta ahora usaste imágenes que hizo otro: \texttt{hello-world} y \texttt{nginx}. Hoy verás cómo se arma una \textbf{propia}, con la API de Café Nube: así corre igual en tu laptop, en la de Diego o en un servidor, \textbf{sin instalar Python ni Flask}. Y cuando una app tiene varias piezas, \textbf{Compose} las levanta todas con un comando.}
\end{minipage}
\end{center}
\end{frame}

% 7
\separador{Docker desde cero y Compose}{Dockerfile, puertos, variables de entorno y varias piezas}

% 8
\begin{frame}[fragile]{Dockerfile: tu propia imagen}
\begin{block}{Dockerfile}
\small Un archivo de texto, llamado exactamente \texttt{Dockerfile}, con \textbf{los pasos para armar una imagen}. Cada línea empieza con una \textbf{instrucción} en mayúsculas.
\end{block}
\begin{columns}[T,onlytextwidth]
\column{.40\textwidth}
{\scriptsize\bfseries Proyecto \texttt{hola-docker}: Dockerfile}
\begin{lstlisting}
FROM python:3.12-slim
WORKDIR /app
COPY app.py .
CMD ["python", "app.py"]
\end{lstlisting}
{\scriptsize\bfseries app.py}
\begin{lstlisting}[style=python]
import platform
print("Hola desde un contenedor")
print("Este programa corre en:", platform.system())
\end{lstlisting}
\column{.56\textwidth}
\renewcommand{\arraystretch}{1.3}
\scriptsize
\begin{tabular}{@{}>{\ttfamily\bfseries\color{utpRojo}}l>{\raggedright\arraybackslash}p{5.6cm}@{}}
\toprule
\textsf{\textbf{\color{black}Instrucción}} & \textbf{Qué hace}\\
\midrule
FROM & Elige la \textbf{imagen base}, de la que partes: \texttt{python:3.12-slim} ya trae Python 3.12\\
WORKDIR & La carpeta, dentro de la imagen, donde pasa todo lo que sigue\\
COPY & \texttt{COPY origen destino}: copia un archivo de tu proyecto a la imagen. El \texttt{.} es la carpeta de \texttt{WORKDIR}\\
CMD & El comando que corre \textbf{cuando arranca} un contenedor: \texttt{python app.py}, escrito como lista\\
\bottomrule
\end{tabular}
\end{columns}
\end{frame}

% 9
\begin{frame}[fragile]{Construir y correr}
\begin{center}
\begin{tikzpicture}[font=\small]
  \node[caja,fill=utpGrisClaro,text width=2.6cm,minimum height=1.2cm] (df) at (0,0) {\textbf{Dockerfile}\\\footnotesize y \texttt{app.py}};
  \node[caja,fill=utpRojo!30,text width=2.2cm,minimum height=1.2cm] (img) at (5.0,0) {\textbf{Imagen}\\\footnotesize\texttt{hola}};
  \draw[flecha] (df) -- node[above,font=\scriptsize\ttfamily,align=center]{docker build\\-t hola .} node[below,font=\scriptsize]{una vez} (img);
  \foreach \n/\y in {1/0.6,2/-0.6} {
    \node[caja,draw=okVerde,fill=okVerde!12,minimum width=2.5cm,minimum height=0.6cm,font=\scriptsize] (k\n) at (10.2,\y) {contenedor \n};
    \draw[flecha] (img.east) -- (k\n.west);
  }
  \node[font=\scriptsize\ttfamily] at (10.2,1.2) {docker run hola};
  \node[font=\scriptsize] at (10.2,-1.2) {todas las veces que quieras};
\end{tikzpicture}
\end{center}
\begin{columns}[c,onlytextwidth]
\column{.54\textwidth}
\begin{lstlisting}[style=terminal]
$ docker build -t hola .
 => [1/3] FROM docker.io/library/python:3.12-slim
 => [2/3] WORKDIR /app
 => [3/3] COPY app.py .
 => naming to docker.io/library/hola
$ docker run hola
Hola desde un contenedor
Este programa corre en: Linux
\end{lstlisting}
\column{.42\textwidth}
\begin{block}{Construir (\textit{build})}
\small Armar una imagen siguiendo el Dockerfile \textbf{de arriba hacia abajo}. Aunque uses Windows o Mac, dentro del contenedor es \textbf{Linux}.
\end{block}
\begin{alertblock}{Ojo}
\small \texttt{CMD} no corre al construir: corre al arrancar. Y si cambias \texttt{app.py}, \textbf{construye otra vez}.
\end{alertblock}
\end{columns}
\end{frame}

% 10
\begin{frame}[fragile]{Puertos: publicar un puerto}
\begin{center}
\begin{tikzpicture}[font=\small]
  \node[zona,minimum width=13.4cm,minimum height=2.5cm] (pc) at (5.5,0) {};
  \node[font=\scriptsize\bfseries,anchor=north west] at (pc.north west) {\ Tu computadora};
  \node[caja,fill=utpGrisClaro,text width=2.4cm,minimum height=1.0cm] (nav) at (0.9,-0.15) {\textbf{Navegador}\\\footnotesize\texttt{localhost:8080}};
  \node[caja,fill=utpRojo!30,text width=1.6cm,minimum height=1.0cm] (p) at (5.0,-0.15) {\textbf{8080}\\\footnotesize tu puerto};
  \node[caja,dashed,draw=okVerde,minimum width=3.6cm,minimum height=1.7cm] (k) at (10.0,-0.15) {};
  \node[font=\scriptsize\bfseries,anchor=north] at (k.north) {Contenedor \texttt{web}};
  \node[caja,draw=okVerde,fill=okVerde!12,text width=2.4cm,minimum height=0.8cm] (ng) at (10.0,-0.35) {\textbf{Nginx}\\\footnotesize escucha en el \textbf{80}};
  \draw[flecha] (nav) -- node[above,font=\scriptsize]{pedido} (p);
  \draw[flecha,utpRojo] (p) -- node[above,font=\scriptsize\ttfamily]{-p 8080:80} (ng);
\end{tikzpicture}
\end{center}
\vspace{1mm}
\begin{columns}[T,onlytextwidth]
\column{.48\textwidth}
\begin{block}{Puerto}
\small Un número que dice \textbf{a qué programa} de una computadora va un pedido. Nginx escucha en el 80; la API de Café Nube, en el 8080.
\end{block}
\begin{block}{Publicar un puerto}
\small \texttt{-p tu-puerto:puerto-del-contenedor}. El contenedor está \textbf{aislado}: sin \texttt{-p}, nada de afuera le llega.
\end{block}
\column{.48\textwidth}
\begin{lstlisting}
docker run -d --name web -p 8080:80 nginx
curl localhost:8080
docker rm -f web   # detiene y borra
\end{lstlisting}
\begin{alertblock}{port is already allocated}
\small Dos contenedores no pueden usar el mismo puerto de \textbf{tu computadora}. El del contenedor sí se repite: \texttt{-p 8081:80}.
\end{alertblock}
\end{columns}
\end{frame}

% 11
\begin{frame}[fragile]{Variables de entorno}
\begin{center}
\begin{tikzpicture}[font=\small]
  \node[caja,fill=utpRojo!30,text width=2.4cm,minimum height=1.4cm] (img) at (0,0) {\textbf{Imagen}\\\footnotesize\texttt{cafe-api}\\\footnotesize una sola};
  \node[caja,draw=okVerde,fill=okVerde!12,text width=4.2cm,minimum height=0.8cm,align=left,font=\scriptsize] (a) at (7.2,0.75) {\textbf{api}: sin variable\\API en el \textbf{8080} · \texttt{-p 8080:8080}};
  \node[caja,draw=okVerde,fill=okVerde!12,text width=4.2cm,minimum height=0.8cm,align=left,font=\scriptsize] (b) at (7.2,-0.75) {\textbf{api-2}: \texttt{-e PUERTO=9090}\\API en el \textbf{9090} · \texttt{-p 8081:9090}};
  \draw[flecha] (img.east) -- (a.west);
  \draw[flecha] (img.east) -- (b.west);
\end{tikzpicture}
\end{center}
\vspace{2mm}
\begin{columns}[T,onlytextwidth]
\column{.48\textwidth}
\begin{block}{Variable de entorno}
\small Un valor con nombre, como \texttt{PUERTO=9090}, que el programa lee \textbf{al arrancar}. Se le pasa al contenedor con \texttt{-e NOMBRE=valor}: un \texttt{-e} por cada una.
\end{block}
\column{.48\textwidth}
\begin{block}{Para qué sirve}
\small Lo que cambia de un lugar a otro (un puerto, una contraseña, el nombre de la tienda) no va dentro de la imagen: \textbf{la misma imagen} sirve en todos lados.
\end{block}
\end{columns}
\par\vspace{2mm}
\begin{exampleblock}{Producción}
\small El lugar donde corre la app que usan los clientes de verdad. Usa la misma imagen que tu computadora, con su propia configuración.
\end{exampleblock}
\end{frame}

% 12
\begin{frame}[fragile]{Una base de datos que pide contraseña}
\small Café Nube va a guardar los pedidos en una \textbf{base de datos}: un programa que guarda información ordenada y la entrega cuando se la piden. \textbf{PostgreSQL} tiene su imagen oficial, \texttt{postgres}, y \textbf{no arranca sin contraseña}.
\par\vspace{3mm}
\begin{columns}[T,onlytextwidth]
\column{.49\textwidth}
{\scriptsize\bfseries \mal\ Sin la variable}
\begin{lstlisting}[style=terminal]
$ docker run -d --name db postgres:17
$ docker logs db
Error: Database is uninitialized and
  superuser password is not specified.
  You must specify POSTGRES_PASSWORD
  to a non-empty value for the superuser.
\end{lstlisting}
{\footnotesize El contenedor se detiene enseguida: sus \textbf{logs} dicen por qué.}
\column{.49\textwidth}
{\scriptsize\bfseries \bien\ Con la variable}
\begin{lstlisting}[style=terminal]
$ docker run -d --name db \
    -e POSTGRES_PASSWORD=secreto postgres:17
$ docker logs db
... database system is ready to accept
    connections
$ docker rm -f db
\end{lstlisting}
{\footnotesize \texttt{:17} elige la versión. PostgreSQL escucha en el puerto \textbf{5432}.}
\end{columns}
\par\vspace{3mm}
\begin{alertblock}{Cuidado}
\small Una contraseña real no se escribe en el Dockerfile ni en el código: va en una variable de entorno, al correr el contenedor.
\end{alertblock}
\end{frame}

% 13
\begin{frame}[fragile]{Café Nube: la API en una imagen}
\begin{block}{Enunciado}
\small La API de la carta solo corre en la laptop de Lucía. Hay que empaquetarla para que corra igual en cualquier computadora: debe \textbf{partir de Python 3.12}, \textbf{instalar las librerías de \texttt{requirements.txt}}, \textbf{responder en el puerto 8080} y quedar funcionando \textbf{sin pasos extra}.
\end{block}
\begin{columns}[T,onlytextwidth]
\column{.42\textwidth}
\begin{lstlisting}
FROM python:3.12-slim
WORKDIR /app
COPY requirements.txt .
RUN pip install -r requirements.txt
COPY . .
EXPOSE 8080
CMD ["python", "api.py"]
\end{lstlisting}
\column{.55\textwidth}
\renewcommand{\arraystretch}{1.2}
\scriptsize
\begin{tabular}{@{}>{\raggedright\arraybackslash}p{2.6cm}>{\raggedright\arraybackslash}p{3.9cm}@{}}
\toprule
\textbf{El enunciado dice\dots} & \textbf{\dots y se vuelve}\\
\midrule
partir de Python 3.12 & \texttt{FROM python:3.12-slim}\\
instalar las librerías & \texttt{COPY requirements.txt} y \texttt{RUN pip install}\\
(el código de la API) & \texttt{COPY . .}\\
responder en el 8080 & \texttt{EXPOSE 8080}, y al correrla \texttt{-p 8080:8080}\\
sin pasos extra & \texttt{CMD [\string"python\string", \string"api.py\string"]}\\
\bottomrule
\end{tabular}
\end{columns}
\par\vspace{2mm}
\begin{columns}[T,onlytextwidth]
\column{.48\textwidth}
{\small\textbf{\color{utpRojo}RUN}: ejecuta un comando \textbf{mientras se construye} la imagen. Aquí instala Flask.}
\column{.48\textwidth}
{\small\textbf{\color{utpRojo}EXPOSE}: anota en qué puerto escucha la app. Solo lo documenta: el puerto lo abre \texttt{-p}.}
\end{columns}
\end{frame}

% 14
\begin{frame}{Capas: por qué importa el orden}
\begin{center}
\begin{tikzpicture}[font=\small]
  \node[font=\footnotesize\bfseries] at (2.8,2.7) {Primera vez};
  \node[font=\footnotesize\bfseries] at (9.2,2.7) {Cambias \texttt{api.py} y construyes};
  \node[capa,fill=utpRojo!20] at (0,2.05) {[1/5] FROM python:3.12-slim};
  \node[capa,fill=utpRojo!20] at (0,1.50) {[2/5] WORKDIR /app};
  \node[capa,fill=utpRojo!20] at (0,0.95) {[3/5] COPY requirements.txt .};
  \node[capa,fill=utpRojo!20] at (0,0.40) {[4/5] RUN pip install -r ...};
  \node[capa,fill=utpRojo!20] at (0,-0.15) {[5/5] COPY . .};
  \node[capa,draw=utpGrisMedio,fill=utpGrisClaro,text=utpGris] at (6.4,2.05) {CACHED [1/5] FROM python:3.12-slim};
  \node[capa,draw=utpGrisMedio,fill=utpGrisClaro,text=utpGris] at (6.4,1.50) {CACHED [2/5] WORKDIR /app};
  \node[capa,draw=utpGrisMedio,fill=utpGrisClaro,text=utpGris] at (6.4,0.95) {CACHED [3/5] COPY requirements.txt .};
  \node[capa,draw=utpGrisMedio,fill=utpGrisClaro,text=utpGris] at (6.4,0.40) {CACHED [4/5] RUN pip install -r ...};
  \node[capa,fill=utpRojo!20] at (6.4,-0.15) {[5/5] COPY . .};
  \node[font=\scriptsize,anchor=west,text=utpGris] at (12.1,0.95) {se reutilizan};
  \node[font=\scriptsize\bfseries,anchor=west,text=utpRojo] at (12.1,-0.15) {se rehace};
\end{tikzpicture}
\end{center}
\begin{block}{Capa}
\small Lo que deja cada línea del Dockerfile, guardado aparte. Si una línea y las de arriba no cambiaron, Docker \textbf{reutiliza su capa} (\texttt{CACHED}) en vez de rehacerla.
\end{block}
\begin{exampleblock}{Por eso, este orden}
\small Primero \texttt{requirements.txt} y \texttt{pip install}; \textbf{después} \texttt{COPY . .}. Así, un cambio en \texttt{api.py} no vuelve a instalar Flask.
\end{exampleblock}
\end{frame}

% 15
\begin{frame}{Café Nube: errores frecuentes}
\renewcommand{\arraystretch}{1.45}
\centering\small
\begin{tabular}{@{}>{\raggedright\arraybackslash}p{5.4cm}>{\raggedright\arraybackslash}p{7.2cm}@{}}
\toprule
\textbf{Lo que ves} & \textbf{Por qué y cómo se arregla}\\
\midrule
\texttt{No module named \textquotesingle flask\textquotesingle} en \texttt{docker logs api} & Faltó instalar las librerías: agrega \texttt{RUN pip install -r requirements.txt} y construye otra vez\\
\texttt{curl} no conecta & Faltó \texttt{-p 8080:8080} al correrla, o el contenedor se detuvo: mira \texttt{docker logs api}\\
\texttt{The container name \string"/api\string" is already in use} & Ya hay un contenedor \texttt{api}: bórralo con \texttt{docker rm -f api}\\
\texttt{port is already allocated} & Otro contenedor usa ese puerto de tu computadora: elige otro, como \texttt{-p 8081:8080}\\
El cambio en \texttt{api.py} no aparece & La imagen tiene la copia vieja: \texttt{docker build -t cafe-api .} otra vez\\
\bottomrule
\end{tabular}
\par\vspace{4mm}
\begin{alertblock}{Lo primero que miras}
\small Cuando un contenedor falla o se detiene solo, \texttt{docker logs} dice por qué.
\end{alertblock}
\end{frame}

% 16
\begin{frame}{Una aplicación, varias piezas}
\begin{columns}[c,onlytextwidth]
\column{.42\textwidth}
\centering
\begin{tikzpicture}[font=\small]
  \node[caja,dashed,text width=3.6cm,minimum height=0.6cm] (nav) at (0,2.9) {Navegador del cliente};
  \node[caja,fill=utpGrisClaro,text width=3.6cm,minimum height=0.9cm] (web) at (0,1.5) {\textbf{web}: la página\\\footnotesize frontend};
  \node[caja,fill=utpRojo!30,text width=3.6cm,minimum height=0.9cm] (api) at (0,0) {\textbf{api}: la API de la carta\\\footnotesize backend};
  \node[caja,draw=okVerde,fill=okVerde!12,text width=3.6cm,minimum height=0.9cm] (db) at (0,-1.5) {\textbf{db}: PostgreSQL\\\footnotesize los pedidos};
  \draw[flecha] (nav) -- node[right,font=\scriptsize]{abre la página} (web);
  \draw[flecha] (web) -- node[right,font=\scriptsize]{pide la carta} (api);
  \draw[flecha] (api) -- node[right,font=\scriptsize]{guarda y lee} (db);
\end{tikzpicture}
\par\vspace{1mm}{\scriptsize Cada pieza, en su propio contenedor.}
\column{.54\textwidth}
\begin{block}{Frontend y backend}
\small El \textbf{frontend} es lo que se ve en el navegador: la página de la carta. El \textbf{backend} es la API que le entrega los datos.
\end{block}
\begin{block}{Servicio}
\small Cada pieza de la app, con su imagen y su configuración. Café Nube tiene \textbf{tres}: \texttt{web}, \texttt{api} y \texttt{db}.
\end{block}
\begin{exampleblock}{Un programa por contenedor}
\small Así cambias la página sin tocar la API, y reinicias la API sin apagar la base.
\end{exampleblock}
\end{columns}
\end{frame}

% 17
\begin{frame}[fragile]{Docker Compose: un archivo, un comando}
\begin{columns}[T,onlytextwidth]
\column{.52\textwidth}
{\small\bfseries\mal\ Con \texttt{docker run}}
\begin{lstlisting}[basicstyle=\ttfamily\scriptsize]
docker build -t cafe-api .
docker run -d --name api -p 8080:8080 cafe-api
docker run -d --name web -p 8081:80 nginx:1.29
\end{lstlisting}
{\footnotesize Un comando largo por pieza, en orden, y hay que recordarlos cada vez. Con la base de datos, uno más.}
\column{.44\textwidth}
{\small\bfseries\bien\ Con Docker Compose}
\begin{lstlisting}[style=yaml]
services:
  api:
    build: .
    ports:
      - "8080:8080"
  web:
    image: nginx:1.29
    ports:
      - "8081:80"
\end{lstlisting}
{\footnotesize Todo en \texttt{compose.yaml}, y un comando: \texttt{docker compose up -d}.}
\end{columns}
\par\vspace{1mm}
\begin{block}{Docker Compose}
\small Una herramienta de Docker que lee los servicios de una app desde un archivo, \texttt{compose.yaml}, y \textbf{los levanta todos con un solo comando}. El archivo va junto al código: quien lo descarga levanta la app igual que tú.
\end{block}
\par\vspace{1mm}
{\footnotesize Viene con Docker Desktop: \texttt{docker compose version}. Se escribe \texttt{docker compose}, \textbf{con espacio}; \texttt{docker-compose}, con guion, es la versión antigua.}
\end{frame}

% 18
\begin{frame}[fragile]{El archivo \texttt{compose.yaml}}
\begin{block}{YAML}
\small Un formato de texto para configuración: \texttt{nombre: valor} en cada línea. \textbf{La sangría dice qué va dentro de qué}, y una lista se escribe con guion: \texttt{- elemento}.
\end{block}
\begin{columns}[c,onlytextwidth]
\column{.28\textwidth}
\begin{lstlisting}[style=yaml]
services:
  api:
    build: .
    ports:
      - "8080:8080"
  web:
    image: nginx:1.29
    ports:
      - "8081:80"
\end{lstlisting}
\column{.68\textwidth}
\renewcommand{\arraystretch}{1.3}
\scriptsize
\begin{tabular}{@{}>{\ttfamily}l>{\raggedright\arraybackslash}p{3.7cm}>{\raggedright\arraybackslash}p{2.6cm}@{}}
\toprule
\textsf{\textbf{Clave}} & \textbf{Qué dice} & \textbf{Con \texttt{docker}}\\
\midrule
services: & La lista de servicios, cada uno con dos espacios de sangría & ---\\
build: . & Construye la imagen con el \texttt{Dockerfile} de esta carpeta & \texttt{docker build .}\\
image: nginx:1.29 & Usa una imagen ya hecha, sin construir nada & \texttt{docker run nginx:1.29}\\
ports: & Los puertos publicados, uno por línea & \texttt{-p 8081:80}\\
\bottomrule
\end{tabular}
\end{columns}
\end{frame}

% 19
\begin{frame}{Compose: up, ps, logs y down}
\begin{center}
\begin{tikzpicture}[font=\small]
  \node[caja,fill=utpGrisClaro,text width=2.6cm,minimum height=1.0cm] (f) at (0,0) {\footnotesize\texttt{compose.yaml}\\\footnotesize los servicios};
  \node[caja,draw=okVerde,fill=okVerde!12,text width=2.6cm,minimum height=1.0cm] (r) at (5.4,0) {\textbf{Corriendo}\\\footnotesize se miran con \texttt{ps} y \texttt{logs}};
  \node[caja,draw=utpGrisMedio,dashed,text width=2.6cm,minimum height=1.0cm] (b) at (10.8,0) {\textbf{Borrados}\\\footnotesize las imágenes quedan};
  \draw[flecha] (f) -- node[above,font=\scriptsize\ttfamily,align=center]{docker compose\\up -d} (r);
  \draw[flecha] (r) -- node[above,font=\scriptsize\ttfamily,align=center]{docker compose\\down} (b);
\end{tikzpicture}
\end{center}
\renewcommand{\arraystretch}{1.15}
\centering\small
\begin{tabular}{@{}>{\ttfamily\footnotesize}l>{\ttfamily\footnotesize}l>{\raggedright\arraybackslash}p{4.4cm}@{}}
\toprule
\textsf{\small\textbf{Con docker}} & \textsf{\small\textbf{Con docker compose}} & \textbf{Qué hace}\\
\midrule
docker run -d, uno por uno & docker compose up -d & Construye lo que falte y levanta todo\\
docker ps & docker compose ps & Lista los contenedores del proyecto\\
docker logs api & docker compose logs api & Los mensajes de un servicio\\
docker rm -f, uno por uno & docker compose down & Detiene y borra todos los contenedores\\
\bottomrule
\end{tabular}
\par\vspace{2mm}
\begin{alertblock}{Cambiaste el código}
\small \texttt{up -d} reutiliza la imagen que ya existe. Para que construya otra vez: \texttt{docker compose up -d -{}-build}.
\end{alertblock}
\end{frame}

% 20
\begin{frame}{Cierre: preguntas para comprobar}
\begin{block}{Responde con tus palabras}
\small
1. \texttt{docker build -t cafe-api .} terminó, pero la API no está encendida. ¿Qué línea del Dockerfile la enciende y cuándo corre?\\\medskip
2. Dos contenedores de Nginx escuchan en su puerto 80. ¿Pueden correr a la vez? ¿Con qué \texttt{-p}?\\\medskip
3. En \texttt{-e PUERTO=9090 -p 8081:9090}, ¿qué número es de tu computadora y cuál del contenedor?\\\medskip
4. Agregas un café a \texttt{api.py}. ¿Qué haces para verlo en el contenedor? ¿Por qué no se vuelve a instalar Flask?\\\medskip
5. Después de \texttt{docker compose down}, ¿qué queda en tu computadora: los contenedores o las imágenes?
\end{block}
\begin{alertblock}{Idea final}
\small El \textbf{Dockerfile} lleva los pasos para armar la imagen; \texttt{-p} abre el camino hasta el contenedor y \texttt{-e} le da su configuración sin cambiar la imagen. Cuando la app tiene varias piezas, \textbf{Compose} las describe en un archivo y las maneja con un comando.
\end{alertblock}
\end{frame}

% 21
\begin{frame}{Referencias}
\scriptsize
\begin{itemize}\setlength\itemsep{4pt}
  \item Ruta \textit{Despliegue con Docker}: curso 1, \textit{Docker desde cero} (Dockerfile, puertos, variables de entorno y el escenario Café Nube), y curso 2, \textit{Varios contenedores con Compose} (lecciones 1 y 2).
  \item Docker, Dockerfile y Docker Compose: \texttt{docs.docker.com}. Imágenes \texttt{python}, \texttt{nginx} y \texttt{postgres}: \texttt{hub.docker.com}.
  \item Material previo del curso: Sesión 27 (desplegar, imagen, contenedor y su ciclo de vida, y la API de Café Nube).
  \item Las respuestas de las preguntas están en \texttt{../respuestas.md}.
\end{itemize}
\end{frame}

\end{document}
